\section{Open Questions and Next Steps}
\label{sec:open-questions}

The following five questions remain genuinely open after this replication. Each has a concrete, CPU-feasible probe.

\begin{enumerate}
  \item \textbf{Gate-budget-normalized comparison.} Does the DC-QAOA advantage at $p=1$ survive when re-plotted against a fixed two-qubit-gate budget instead of a fixed number of QAOA layers?
  \begin{itemize}
    \item \emph{Basis:} DC-QAOA adds $2|E|$ extra two-qubit rotations per layer ($Z_iY_j$ and $Y_iZ_j$ per edge) on top of QAOA's $|E|$ $ZZ$ rotations, so a ``$p=1$ DC-QAOA'' circuit uses ${\sim}3\times$ the two-qubit gates of ``$p=1$ QAOA''. The paper compares vs $p$, not vs gate count, which flatters the CD variant on NISQ-relevant metrics.
    \item \emph{Next steps:} Re-run the same $n\in\{4,6,8\}$ 3-regular MaxCut sweep, but bin results by total two-qubit-gate count $G$ rather than depth $p$. Compare best-$R$-at-fixed-$G$ for QAOA and DC-QAOA. If DC-QAOA still wins at matched $G$, the advantage is real for NISQ; if it disappears, the paper's plot choice is misleading. CPU-only, ${\sim}1$~hr.
  \end{itemize}

  \item \textbf{Adaptive CD vs fixed pool.} Does the CD operator pool $\{Z\otimes Y, Y\otimes Z\}$ used in this paper remain optimal, or does an ADAPT-QAOA / adaptive-CD selection procedure that grows the pool problem-by-problem beat it at matched depth?
  \begin{itemize}
    \item \emph{Basis:} The paper fixes a hand-picked two-body CD pool derived from Nested Commutator expansion of the adiabatic-gauge potential. ADAPT-QAOA (Zhu \emph{et al.} 2022) instead grows the pool one operator at a time chosen by gradient magnitude; adaptive-CD (Hegade \emph{et al.} 2022) does similar for CD terms. The claim of DC-QAOA's low-$p$ optimality is not conditional on ``the best available adaptive scheme''.
    \item \emph{Next steps:} Implement a small adaptive-CD variant that picks each layer's CD term from a candidate pool ($\{ZY, YZ, XY, YX, ZX, XZ\}$) by gradient magnitude on the current state. Compare best-$R$ at $p=1..3$ on the same $n=6, n=8$ 3-regular MaxCut instances. If adaptive-CD equals or beats fixed DC-QAOA at matched $p$, the paper's pool choice is a hyperparameter not a discovery.
  \end{itemize}

  \item \textbf{Noise robustness of CD terms.} How robust is the DC-QAOA advantage to realistic gate noise, given the extra two-qubit gates?
  \begin{itemize}
    \item \emph{Basis:} The paper motivates DC-QAOA as an NISQ-friendly shallow-circuit method, but presents only noiseless statevector results (as does this replication). The extra $Z_iY_j$ and $Y_iZ_j$ gates each require 2 CNOTs plus single-qubit rotations under standard decomposition; on today's ${\sim}99.5\%$ two-qubit-fidelity hardware, one $p=1$ DC-QAOA layer accumulates roughly $3\times$ the two-qubit error of one QAOA layer.
    \item \emph{Next steps:} Use \texttt{qiskit-aer}'s depolarizing noise model at per-two-qubit-gate error rates $\{0.001, 0.005, 0.01\}$ on the same $K_4$ and $n=6$ 3-regular instances at $p=1,2$. Plot $R$ vs noise strength for both variants; identify the crossover error rate below which DC-QAOA still wins. CPU-feasible in $<1$~hr per instance.
  \end{itemize}

  \item \textbf{Generalization to weighted / realistic problem classes.} Does the DC-QAOA low-$p$ advantage extend to weighted / more-realistic problem classes (weighted MaxCut, portfolio optimization, MIS on non-regular graphs) where the QAOA landscape structure differs?
  \begin{itemize}
    \item \emph{Basis:} The paper tests only unweighted 3-regular MaxCut, LFIM, Sherrington--Kirkpatrick and $p$-spin --- all with high symmetry or Gaussian-random couplings. Real-world combinatorial problems (Markowitz portfolio, weighted MaxCut from graph clustering) have skewed weight distributions and no permutation symmetry, which can dominate the QAOA landscape structure.
    \item \emph{Next steps:} Repeat the DC-QAOA vs QAOA sweep on: (a) weighted MaxCut with edge weights drawn from $|\mathcal{N}(0,1)|$ on $n=6,8$ graphs; (b) small ($n=6$) Markowitz portfolio-optimization QUBO built from 6 real S\&P sector returns. Look for whether DC-QAOA's $p=1$ advantage generalizes or is specific to the paper's symmetric instances.
  \end{itemize}

  \item \textbf{ML-optimized CD schedules.} Can machine-learning-optimized CD schedules (rather than a single scalar $\alpha$ per layer) further compress QAOA depth while preserving the DC-QAOA advantage?
  \begin{itemize}
    \item \emph{Basis:} The paper uses a single variational parameter $\alpha_k$ per DC layer, giving a piecewise-constant CD schedule. Recent work on shortcuts-to-adiabaticity uses ML-optimized time-dependent CD amplitudes (Sels--Polkovnikov 2017 extended by Gu\'ery-Odelin 2019) that can further reduce time-to-adiabatic-fidelity by 3--10$\times$ in some models. Whether this transfers to the digitized/variational setting is untested.
    \item \emph{Next steps:} Parametrize the CD amplitude within each DC layer as a small MLP $\alpha(\tau;\theta)$ sampled at $K$ intermediate Trotter substeps ($K=3, 5$), then variationally optimize $\theta$ jointly with QAOA $\beta,\gamma$. Compare to the fixed-$\alpha$ DC-QAOA baseline on $n=6$ 3-regular MaxCut at $p=1,2$. If the ML-CD variant reaches higher $R$ at matched two-qubit-gate count, it opens a path to sub-layer-count QAOA. CPU-feasible; needs ${\sim}50$ restarts due to larger param count.
  \end{itemize}
\end{enumerate}
